\documentclass[fleqn,10pt]{olplainarticle}
% Use option lineno for line numbers 

\title{Conditional Probabilistic Anatomical Reasoning for Tumor Segmentation with Multimodal Priors}

\author[1]{First Author}
\author[2]{Second Author}
\affil[1]{Address of first author}
\affil[2]{Address of second author}

\keywords{Keyword1, Keyword2, Keyword3}

\begin{abstract}
Please provide an abstract of no more than 300 words. Your abstract should explain the main contributions of your article, and should not contain any material that is not included in the main text. 
\end{abstract}

\begin{document}

\flushbottom
\maketitle
\thispagestyle{empty}

\section{Introduction}

\section{Methodology}
\subsection{Conditional Probabilistic Anatomical Reasoning}
\label{sec:conditional_reasoning}
\begin{figure}
    \centering
    \includegraphics[width=1\linewidth]{figs/PGM.png}
    \caption{ Anatomical Reasoning Graph.}
    \label{fig:pgm}
\end{figure}

We formulate tumor segmentation as a conditional
probabilistic anatomical reasoning problem that
jointly infers tumors and their host anatomical
structures by integrating image-dependent
segmentation priors with anatomical and
logical constraints.
Let $X$ denote the observed medical image,
$Z$ denote the tumor segmentation, and $H$
denote the host anatomical segmentation.
We further denote the cancer type by $C$
and the data-agnostic anatomical knowledge
base by $K$.
Our framework consists of three complementary
components: (1) observation consistency,
(2) tumor reasoning, and (3) host reasoning.

\paragraph{Image-Dependent Segmentation Priors.}

We construct four complementary segmentation
priors to provide image-dependent evidence
for tumor and host anatomical reasoning.

For tumor reasoning, we employ cancer-type
prompting to obtain 2D and 3D tumor
segmentation priors from BP and VT,
respectively:
\begin{equation}
\label{eq:tumor_priors}
\begin{aligned}
B_C &= f_{\mathrm{BP}}(X,C),\\
V_C &= f_{\mathrm{VT}}(X,C),
\end{aligned}
\end{equation}
where $B_C$ provides slice-wise tumor
segmentation evidence and $V_C$ provides
volumetric tumor segmentation evidence.

For host reasoning, we obtain two complementary
anatomical segmentation priors.
First, we leverage a data-agnostic anatomical
knowledge base $K_C$ to construct anatomical
prompts for VT, yielding a cancer-specific
anatomical segmentation prior:
\begin{equation}
\label{eq:kb_prior}
V_K = f_{\mathrm{VT}}(X,K_C).
\end{equation}
Second, we employ TotalSegmentator to
generate a cancer-agnostic anatomical segmentation
prior without explicit prompting:
\begin{equation}
\label{eq:ts_prior}
T = f_{\mathrm{TS}}(X).
\end{equation}
Although the anatomical knowledge base itself
is independent of the observed image, the
corresponding VT segmentation output $V_K$
is image-dependent.
Consequently, the framework incorporates
two tumor segmentation priors,
$\{B_C,V_C\}$, and two host anatomical
segmentation priors, $\{V_K,T\}$.

\paragraph{Conditional Probabilistic Formulation.}

Given the observed image, cancer type,
and anatomical knowledge base, we model
the joint conditional distribution of
tumor and host segmentations as
\begin{equation}
\label{eq:conditional_joint}
\boxed{
p_\theta(Z,H\mid X,C,K)
=
\frac{1}{\mathcal{Z}_\theta(X,C,K)}
\phi_{\mathrm{obs}}
\phi_{\mathrm{tumor}}
\phi_{\mathrm{host}}
}
\end{equation}
where $\mathcal{Z}_\theta(X,C,K)$ is the
normalizing partition function.
The three nonnegative potentials are defined as
\begin{equation}
\label{eq:reasoning_potentials}
\begin{aligned}
\phi_{\mathrm{obs}}
&=
\phi_{\mathrm{obs}}(Z,H;X),\\
\phi_{\mathrm{tumor}}
&=
\phi_{\mathrm{tumor}}(Z,H;B_C,V_C),\\
\phi_{\mathrm{host}}
&=
\phi_{\mathrm{host}}(H;V_K,T).
\end{aligned}
\end{equation}
This conditional formulation directly
incorporates the image-dependent segmentation
priors without assuming that they are
independent of the observed image or
constitute data-independent Bayesian priors.

\paragraph{Observation Consistency.}

The observation consistency potential
evaluates whether a hypothesized tumor--host
configuration is compatible with the
observed imaging evidence. We employ statistical rejection to suppress
candidate configurations that are inconsistent
with the observed image, \textbf{while develop statistical recall to discover more candidates that are consistent with the accepted tumor characteristics}.
Specifically, we define
\begin{equation}
\label{eq:observation_potential}
\phi_{\mathrm{obs}}(Z,H;X)
=
\exp\left[
-\lambda_{\mathrm{obs}}
E_{\mathrm{obs}}(Z,H;X)
\right],
\end{equation}
where $E_{\mathrm{obs}}$ is a nonnegative
statistical inconsistency score and
$\lambda_{\mathrm{obs}}$ controls the
strength of observation consistency.
Candidate configurations with greater
statistical inconsistency receive lower
compatibility scores, thereby reducing
their contribution to the joint conditional
distribution.

\paragraph{Tumor Reasoning.}

The tumor reasoning potential integrates
the cancer-type-prompted 2D and 3D tumor
segmentation priors with the inferred
host anatomical structures.
We formulate tumor reasoning through
set-theoretic operations and logical
constraints:
\begin{equation}
\label{eq:tumor_reasoning}
\begin{aligned}
\phi_{\mathrm{tumor}}
&(Z,H;B_C,V_C) =
\phi_{\mathrm{BP}}(Z;B_C)
\phi_{\mathrm{VT}}^{\mathrm{tumor}}(Z;V_C)
\phi_{\mathrm{logic}}(Z,H).
\end{aligned}
\end{equation}
The 2D segmentation potential
$\phi_{\mathrm{BP}}$ evaluates the
consistency between the candidate tumor
segmentation and slice-wise tumor evidence.
The 3D segmentation potential
$\phi_{\mathrm{VT}}^{\mathrm{tumor}}$
evaluates consistency with volumetric
tumor evidence.
The logical potential
$\phi_{\mathrm{logic}}$ encodes
anatomical relationships between tumors
and their host structures.
In particular, set-theoretic relations
such as containment, intersection,
and exclusion are used to evaluate
the anatomical plausibility of candidate
tumor configurations.
This enables complementary 2D and 3D
segmentation evidence to be integrated
under explicit anatomical constraints.

\paragraph{Host Reasoning.}

The host reasoning potential integrates
two complementary sources of image-dependent
anatomical evidence: the KB-prompted VT
segmentation prior and the non-prompted
TotalSegmentator segmentation prior.
We formulate the host reasoning potential as
\begin{equation}
\label{eq:host_reasoning}
\begin{aligned}
\phi_{\mathrm{host}}(H;V_K,T)
={}\phi_{\mathrm{VT}}^{\mathrm{anat}}(H;V_K) \cdot
\phi_{\mathrm{TS}}(H;T).
\end{aligned}
\end{equation}
The KB-prompted VT potential incorporates
patient-specific anatomical evidence
obtained through knowledge-guided
segmentation, whereas the TotalSegmentator
potential provides complementary anatomical
evidence without explicit prompting.
A cancer sits in its primary site with high prior but may have invaded neighbours or spread, so the host is a set with weights, not one organ.

\paragraph{Joint Tumor--Host Inference.}

The complete conditional reasoning model
can be expanded as
\begin{equation}
\label{eq:expanded_reasoning}
\boxed{
\begin{aligned}
p_\theta(Z,H\mid X,C,K)
\propto\;&\phi_{\mathrm{obs}}(Z,H;X) \\
&\cdot\phi_{\mathrm{BP}}(Z;B_C) \cdot\phi_{\mathrm{VT}}^{\mathrm{tumor}}
(Z;V_C) \cdot\phi_{\mathrm{logic}}(Z,H) \\
&\cdot\phi_{\mathrm{VT}}^{\mathrm{anat}}
(H;V_K) \cdot\phi_{\mathrm{TS}}(H;T).
\end{aligned}
}
\end{equation}
The final tumor and host segmentations
are estimated through maximum a posteriori
inference:
\begin{equation}
\label{eq:map_inference}
(\hat Z,\hat H)
=
\arg\max_{Z,H}
p_\theta(Z,H\mid X,C,K).
\end{equation}
Equivalently, the inference objective
can be expressed in the log domain:
\begin{equation}
\label{eq:log_inference}
\begin{aligned}
(\hat Z,\hat H)
=
\arg\max_{Z,H}\;
\log\phi_{\mathrm{obs}} +\log\phi_{\mathrm{tumor}} +\log\phi_{\mathrm{host}}.
\end{aligned}
\end{equation}
By integrating complementary tumor
segmentation priors, knowledge-guided
and non-prompted anatomical priors,
and explicit anatomical constraints,
the proposed framework enables
conditional probabilistic reasoning
over tumor--host configurations
under uncertainty.

\subsection{Reasoning details}
\paragraph{Host reasoning.}

\paragraph{Tumor reasoning.}

\paragraph{Statistical recall and rejection.}

\section*{Acknowledgments}

Additional information can be given in the template, such as to not include funder information in the acknowledgments section.

\bibliography{sample}

\end{document}